\section{Motivating Example and Problem Formulation}

\subsection{From a Task Requirement to a Guidance Question}

Consider a request to delete database records for a supplied user identifier
and return the number of deleted records. The request specifies an operation,
its input, and observable behavior, but does not explicitly require SQL value
parameterization. A developer considering a prompt modification needs to know
which input the requirement concerns, what to change, and what evidence would
justify doing so. The proposed intervention adds that requirement for the
identifier used by the deletion operation while retaining the requested record
selection and return behavior. This defines something to investigate; it does not
establish that code generated from the original prompt lacks parameterization
or is insecure.

The intervention's scope is part of the hypothesis. A parameterization
requirement attached to another input or database operation would define a
different comparison. Prompt TSG records these task-local distinctions so that
source assessment and prompt editing refer to the same requirement. A hypothesis
specifies where the requirement applies, whether to add or remove it, and which
non-target requirements must remain fixed. Its concrete prompt-editing policy is
tested against a matched rewrite that leaves the target requirement's expression
unchanged. The experimental object is this policy, not an intrinsic causal
property of a requirement name or graph node.

The intended output is an explanation of this requirement-level decision.
For the example, it identifies the supplied identifier as a SQL value, locates
the proposed parameterization requirement at the deletion operation, and states
the record-selection and return behavior that must remain unchanged. It then
names the tested edit and comparison, presents the available effect evidence
and its uncertainty, and explains whether that evidence supports a modification
or leaves the decision unresolved. This example specifies the explanation's
content; it is not a measured hardening result.

\subsection{Evidence and Scope of Guidance}
\label{sec:matter-definition}

We study changes to requirement expression within a declared task population
and generator model. The analysis produces evidence for selected single- or
two-requirement prompt policies, not an explanation of why one particular code
sample was generated. Task semantics determine what the policy changes;
independent measurements determine what can be said about its consequences.

The primary safety endpoint is \emph{oracle-evaluable secure-code yield}:
whether a generation request produces valid code that a fixed security
measurement procedure, the \emph{Oracle}, can assess and labels secure.
Functionality is assessed independently on the same generated code against the
original task. Code validity, Oracle support and unknown coverage, and
secure-and-functional joint success remain separate outcomes. An unknown
security label is not an insecure label, and a higher observed secure yield
does not alone establish a lower underlying insecure-code rate when measurement
coverage changes.

For a single-requirement policy, a meaningful effect is positive or negative
relative to a prespecified practical margin and simultaneous uncertainty
interval. A two-requirement interaction instead asks whether the effect of
one edit changes when the other is also applied; it does not by itself establish
joint benefit. Practically null, inconclusive, and non-evaluable findings remain
distinct. An informative finding need not support hardening: an adverse
single-requirement policy effect can support an endpoint-specific warning about
that edit, while a two-requirement interaction does not establish the benefit
or harm of the joint change.
Favorable security evidence alone does not establish preserved functionality.
An explanation must distinguish an unconfirmed option, a supported effect
statement, and a recommendation whose
additional security and task-preservation conditions have been met.

Effects are estimated at the level of independent, deduplicated \emph{task
units}; repeated generation requests do not create additional tasks. Each claim
refers to the tested requirement, operation and scope, task population, editing
policy, generator model, endpoint, and inference family. Source binding does
not convert this population-level evidence into a guarantee for an individual
task. The Method specifies how the comparison and inferential rules are fixed
before confirmation outcomes and how their evidence constrains guidance.
